\section*{ID6940: Simple Pump–Loss Steady-State (Linear-Loss Model): H F ≈ 2.55×10\textasciicircum{}-19 J}
\paragraph{Metadata.}
PiID: 2737231989875 | RTEID: RTE:P26949.5612:T0.4223 | UUID: d76eb43d-d517-5733-a20d-647158fc3366

\paragraph{Canonical Statement.}
Using \textbackslash{}(h f \textbackslash{}approx 2.55\textbackslash{}times10\textasciicircum{}\{-19\}\textbackslash{} \textbackslash{}mathrm\{J\}\textbackslash{}) and \textbackslash{}(\textbackslash{}tau \textbackslash{}approx 2.07\textbackslash{}times10\textasciicircum{}\{-9\}\textbackslash{} \textbackslash{}mathrm\{s\}\textbackslash{}):

\paragraph{Canonical Equation.}
\[
h f \approx 2.55\times10^{-19}\ \mathrm{J}
\]

\paragraph{Dictionary.}
Simple Pump–Loss Steady-State (Linear-Loss Model): H F ≈ 2.55×10\textasciicircum{}-19 J — h f ≈ 2.55×10\textasciicircum{}-19 J

\paragraph{Encyclopedia.}
Simple Pump–Loss Steady-State (Linear-Loss Model): H F ≈ 2.55×10\textasciicircum{}-19 J is a scaling / order-of-magnitude relation, written h f ≈ 2.55×10\textasciicircum{}-19 J. It expresses a scaling or proportionality, capturing how the quantity grows with its parameters. It is built from h, f. Within the Trawin Zero Point registry it serves as one element of the framework's mathematical narrative.
